% Appendix B: Theoretical Analysis of Cross-Attention Design
% For MoM-Driver NeurIPS 2026 Paper

\section{Theoretical Analysis of Cross-Attention Design}
\label{appendix:cross_attention}

This appendix provides a detailed theoretical analysis of our State-Frozen cross-attention mechanism compared to the Linear Interpolation Cross-Attention (LICA) baseline used in prior work~\cite{lady}. While this analysis demonstrates our theoretical understanding, we emphasize that it is \textbf{not claimed as a main contribution} of this paper.

\subsection{Background: Linear Cross-Attention Formulations}

In linear attention architectures, cross-attention between query tokens $\mathbf{Q} \in \mathbb{R}^{N_q \times d}$ and key-value tokens $\{\mathbf{K}, \mathbf{V}\} \in \mathbb{R}^{N_{kv} \times d}$ can be implemented through different approaches. The standard Transformer cross-attention computes:

\begin{equation}
\mathbf{O} = \text{Softmax}\left(\frac{\mathbf{Q}\mathbf{K}^\top}{\sqrt{d}}\right) \mathbf{V}
\label{eq:standard_cross_attn}
\end{equation}

For linear attention systems that maintain recurrent states, cross-attention requires alternative formulations due to the absence of explicit attention matrices.

\subsection{Linear Interpolation Cross-Attention (LICA)}

The LICA approach, used in LADY~\cite{lady} and other multi-modal linear attention models~\cite{vl_mamba, cobra, visual_rwkv}, implements cross-attention by:

\begin{enumerate}
    \item \textbf{Concatenation}: Combine key-value and query sequences: $\mathbf{X} = [\mathbf{K}; \mathbf{V}; \mathbf{Q}] \in \mathbb{R}^{(N_{kv} + N_q) \times d}$
    \item \textbf{Linear Attention}: Apply linear attention (e.g., RWKV, MoM) to the concatenated sequence
    \item \textbf{Extraction}: Extract the query portion: $\mathbf{O} = \text{LinearAttn}(\mathbf{X})[N_{kv}:N_{kv}+N_q, :]$
\end{enumerate}

For RWKV-style linear attention with state $\mathbf{S}_t \in \mathbb{R}^{d \times d}$, this becomes:
\begin{align}
\mathbf{S}_t &= \lambda \mathbf{S}_{t-1} + \mathbf{k}_t \mathbf{v}_t^\top \\
\mathbf{o}_t &= \mathbf{q}_t \mathbf{S}_t
\end{align}

where $\lambda \in (0,1)$ is a decay factor, and $\{\mathbf{k}_t, \mathbf{v}_t, \mathbf{q}_t\}$ are projections of input token $\mathbf{x}_t$.

\subsection{Theoretical Biases in LICA}

We identify two fundamental biases introduced by the LICA formulation:

\subsubsection{Bias 1: Positional Decay Attenuation ($\varepsilon_{decay}$)}

\begin{theorem}[Positional Decay Bias]
In LICA, visual tokens experience artificial attenuation due to their position in the concatenated sequence. For a visual token at position $i$ and query token at position $j = N_{kv} + \delta$ (where $\delta \geq 0$), the effective attention weight is scaled by $\lambda^{\delta}$ compared to standard cross-attention.
\end{theorem}

\begin{proof}
Consider a visual key-value pair $(\mathbf{k}_i, \mathbf{v}_i)$ processed at position $i < N_{kv}$ and a query $\mathbf{q}_j$ at position $j = N_{kv} + \delta$.

In the recurrent state update, when processing token $j$:
\begin{align}
\mathbf{S}_j &= \lambda^{\delta} \mathbf{S}_i + \sum_{t=i+1}^{j-1} \lambda^{j-t-1} \mathbf{k}_t \mathbf{v}_t^\top + \mathbf{k}_j \mathbf{v}_j^\top
\end{align}

The contribution of the visual token $(\mathbf{k}_i, \mathbf{v}_i)$ to the query output is:
\begin{align}
\mathbf{o}_j^{(i)} &= \mathbf{q}_j (\lambda^{\delta} \mathbf{k}_i \mathbf{v}_i^\top) = \lambda^{\delta} (\mathbf{q}_j \mathbf{k}_i) \mathbf{v}_i
\end{align}

Thus, the visual information is attenuated by factor $\lambda^{\delta}$, where $\delta$ is the positional distance. For typical $\lambda \approx 0.9$ and $\delta = N_q$ (number of query tokens), this introduces significant decay:
\begin{equation}
\varepsilon_{decay} = 1 - \lambda^{N_q}
\end{equation}

For $N_q = 64$ (typical trajectory tokens) and $\lambda = 0.9$: $\varepsilon_{decay} \approx 0.9998$, meaning visual information is almost entirely suppressed.
\end{proof}

\subsubsection{Bias 2: Query State Contamination ($\varepsilon_{contam}$)}

\begin{theorem}[Query Contamination Bias]
In LICA, query tokens contribute to the recurrent state, causing later queries to read state contaminated by earlier queries rather than pure visual information.
\end{theorem}

\begin{proof}
When processing query token $\mathbf{q}_j$ at position $j = N_{kv} + \delta$, the state update includes the query's own key-value contribution:
\begin{align}
\mathbf{S}_j &= \lambda \mathbf{S}_{j-1} + \mathbf{k}_j \mathbf{v}_j^\top
\end{align}

For subsequent query $\mathbf{q}_{j+1}$, the output includes contamination from $\mathbf{q}_j$:
\begin{align}
\mathbf{o}_{j+1} &= \mathbf{q}_{j+1} \mathbf{S}_{j+1} \\
&= \mathbf{q}_{j+1} (\lambda \mathbf{S}_j + \mathbf{k}_{j+1} \mathbf{v}_{j+1}^\top) \\
&= \lambda \mathbf{q}_{j+1} \mathbf{S}_j + \mathbf{q}_{j+1} \mathbf{k}_{j+1} \mathbf{v}_{j+1}^\top
\end{align}

The contamination term is:
\begin{equation}
\varepsilon_{contam} = \lambda (\mathbf{q}_{j+1} \mathbf{k}_j) \mathbf{v}_j
\end{equation}

This means query $\mathbf{q}_{j+1}$ reads information written by query $\mathbf{q}_j$, violating the principle that cross-attention should only involve visual tokens.
\end{proof}

\subsection{State-Frozen Cross-Attention}

To eliminate both biases, we propose the State-Frozen mechanism:

\begin{definition}[State-Frozen Cross-Attention]
In State-Frozen cross-attention, visual tokens update the recurrent state normally, but query tokens are prevented from writing to the state. Formally, for positions $t \geq N_{kv}$ (query positions):
\begin{align}
\text{State update:} \quad &\mathbf{S}_t = \mathbf{S}_{t-1} \quad \text{(frozen)} \\
\text{Output computation:} \quad &\mathbf{o}_t = \mathbf{q}_t \mathbf{S}_t \quad \text{(normal)}
\end{align}
\end{definition}

\subsubsection{Bias Elimination}

\begin{theorem}[State-Frozen Bias Elimination]
State-Frozen cross-attention eliminates both $\varepsilon_{decay}$ and $\varepsilon_{contam}$ biases.
\end{theorem}

\begin{proof}
\textbf{Elimination of $\varepsilon_{decay}$:} Since query tokens do not update the state, there is no decay between the end of visual processing and query processing. The state remains $\mathbf{S}_{N_{kv}}$ for all query positions.

\textbf{Elimination of $\varepsilon_{contam}$:} Since query tokens do not write to the state, each query reads the same pure visual state $\mathbf{S}_{N_{kv}}$, with no contamination from other queries.

The output for query $j$ becomes:
\begin{equation}
\mathbf{o}_j = \mathbf{q}_j \mathbf{S}_{N_{kv}} = \mathbf{q}_j \left(\sum_{i=1}^{N_{kv}} \lambda^{N_{kv}-i} \mathbf{k}_i \mathbf{v}_i^\top \right)
\end{equation}

This preserves the standard cross-attention semantics where each query independently accesses visual information.
\end{proof}

\subsection{Empirical Validation}

Despite the theoretical biases, LADY achieves strong performance (PDMS = 90.9) with LICA, suggesting that:

\begin{enumerate}
    \item \textbf{Network Compensation}: Deep networks may learn to compensate for these biases through parameter adaptation
    \item \textbf{Task Tolerance}: Some driving tasks may be robust to attention bias
    \item \textbf{Regularization Effect}: The biases might act as implicit regularization
\end{enumerate}

\subsection{Implementation in MoM-Driver}

In our MoM implementation, State-Frozen is controlled by the \texttt{frozen\_from} parameter:

\begin{verbatim}
# LICA baseline (biased)
frozen_from = -1  # No freezing

# State-Frozen (bias-free)
frozen_from = kv_len  # Freeze from query positions
\end{verbatim}

We compare both modes in our ablations (Section~\ref{sec:ablations}).

\subsection{Complexity Analysis}

\begin{proposition}[Computational Equivalence]
State-Frozen cross-attention has identical computational complexity to LICA: $O(N_{total} \cdot d^2)$ where $N_{total} = N_{kv} + N_q$.
\end{proposition}

The only difference is in the state update logic—the overall computation graph remains unchanged.

\subsection{Connection to Standard Cross-Attention}

\begin{theorem}[Equivalence to Transformer Cross-Attention]
Under certain conditions, State-Frozen cross-attention with infinite precision and no decay ($\lambda = 1$) recovers standard Transformer cross-attention.
\end{theorem}

\begin{proof}
With $\lambda = 1$ and State-Frozen, the output for query $j$ becomes:
\begin{align}
\mathbf{o}_j &= \mathbf{q}_j \left(\sum_{i=1}^{N_{kv}} \mathbf{k}_i \mathbf{v}_i^\top \right) \\
&= \sum_{i=1}^{N_{kv}} (\mathbf{q}_j^\top \mathbf{k}_i) \mathbf{v}_i
\end{align}

This is equivalent to Transformer cross-attention without the softmax normalization, representing a linear attention approximation.
\end{proof}

\subsection{Discussion}

The theoretical analysis reveals that:

\begin{itemize}
    \item \textbf{LICA introduces systematic biases} that can be formally quantified
    \item \textbf{State-Frozen provides a principled solution} while maintaining efficiency
    \item \textbf{Both approaches are viable in practice}, suggesting robustness of deep learning to architectural choices
    \item \textbf{The choice offers a design trade-off} between theoretical purity and empirical optimization
\end{itemize}

While we implement State-Frozen in MoM-Driver for theoretical soundness, we acknowledge that both formulations can be effective, and the choice may depend on specific task requirements and training dynamics.

\subsection{Related Theoretical Work}

Our analysis builds on recent work in linear attention theory:
\begin{itemize}
    \item \textbf{Attention approximation bounds}~\cite{linear_transformer}: Error analysis of linear attention approximations
    \item \textbf{State space theory}~\cite{s4}: Theoretical foundations of recurrent state-based models
    \item \textbf{Memory capacity analysis}~\cite{memory_networks}: Capacity bounds in external memory systems
\end{itemize}

Our contribution lies in identifying and formalizing biases specific to concatenation-based cross-attention in linear models—an important consideration for multi-modal architectures.